\section{Recover a negative reading from its argument}
The argument for a negative integer is one less than its absolute magnitude. In this
project, a negative reading with argument $u$ has value $-1-u$. Thus argument zero gives
minus one, twenty-three gives minus twenty-four, and twenty-four gives minus twenty-five.
There is no integer negative-zero case to add.

\tagpdfsetup{table/header-rows={1}}
\begin{center}
\begin{tabular}{rrl}
Reading & Unsigned argument & Preferred bytes, hexadecimal \\
\hline
23 & 23 & \texttt{17} \\
24 & 24 & \texttt{18 18} \\
$-1$ & 0 & \texttt{20} \\
$-24$ & 23 & \texttt{37} \\
$-25$ & 24 & \texttt{38 18} \\
32767 & 32767 & \texttt{19 7f ff} \\
$-32768$ & 32767 & \texttt{39 7f ff} \\
\end{tabular}
\end{center}

For a required reading, first establish that the major type is unsigned or negative
integer and that its unsigned argument is at most 32767. Then form the signed value
and store its low sixteen bits in the local sample slot. An optional generic negative
integer can have a much larger argument; the walker can validate and skip it without
trying to store its mathematical value in a signed sixty-four-bit register.
